% Provenance (never printed): sources/book2/BRIEFS.md, g12-07. Saving equals investment; private and public saving; net
% lending; the market for loanable funds; crowding out; how firms finance investment; bubbles and crises. Real numbers:
% figdata/grade-12/07-saving-and-investment/ (WDI gross saving and gross capital formation, 2010-2019 means, 155 countries; JST
% USA rgdpmad, unemp, hpnom). Ledger: Feldstein-Horioka 1980 (abstract: the question only; no coefficient read), Federal
% Reserve History on 1929 and 2007-09, Wikipedia on Lehman Brothers. No stock-exchange photograph (Commons rate-limited the download; see PROGRESS). Brenland's
% saving and investment functions are invented.
\chapter{Saving, Investment and the Financial System}\label{ch:g12:saving-and-investment}

From 2010 to 2019, Singapore saved 46\,\% of its national income on average and invested 27\,\% of its GDP at home. Mozambique saved 12\,\% and invested 42\,\%. Across 155
countries, the two rates are related, but loosely: a country that saved ten points more invested, on average, only about two and a half points more.
% figdata: si (SGP 45.7/27.1, MOZ 11.8/42.1); si_key (n 155, slope 0.251, corr 0.477)

In a closed economy, every unit invested must be saved at home. In the world as it is, saving crosses borders, and a country can invest more than it saves by borrowing
from abroad. This chapter starts from the closed economy, where saving equals investment and the interest rate brings them together, then opens it, then looks at the
institutions that move savings to investors, and at what happens when they fail.

\begin{openingchart}{Gross national saving in per cent of gross national income, against gross capital formation (investment) in per cent of GDP, means of 2010--2019,
155 countries; the line marks equal rates.
Three countries with saving below 0 or above 70\,\% lie outside the frame.}{World Bank, World Development Indicators (NY.GNS.ICTR.ZS, NE.GDI.TOTL.ZS)}
\begin{tikzpicture}
\begin{axis}[omscatter, width=0.72\linewidth, height=6.4cm, xmin=0, xmax=70, ymin=0, ymax=70, xlabel={national saving, per cent of national income},
  ylabel={investment, per cent of GDP}, xtick={0,10,...,70}, ytick={0,10,...,70}]
\addplot+[omDef, mark=*, mark size=1.2pt, only marks] table[x=s, y=i] {figdata/grade-12/07-saving-and-investment/out/si.csv};
\addplot[gray, dashed, sharp plot, no marks, domain=0:70, samples=2] {x};
\node[font=\scriptsize, anchor=west] at (axis cs:45.7,27.1) {Singapore};
\node[font=\scriptsize, anchor=west] at (axis cs:11.8,42.1) {Mozambique};
\node[font=\scriptsize, anchor=west] at (axis cs:46.2,44.3) {China};
\end{axis}
\end{tikzpicture}
\end{openingchart}

\begin{recall}
From the grades 6--9 volume: saving, banks, loans and interest. From \cref{ch:g10:national-accounts}: the identity $Y = C + I + G + X - M$. From \cref{ch:g10:price-indices}:
the \omterm{def:g10:price-indices:real-interest-rate}{real interest rate}. From \cref{ch:g11:money-and-banking}: what banks do. From \cref{ch:g12:present-value}: \omterm{def:g12:present-value:bond}{bonds} and their prices.
\end{recall}

\section{Saving equals investment}\label{sec:g12:saving-and-investment:identity}

\begin{definition}[National, private and public saving]\label{def:g12:saving-and-investment:national-saving}
\emph{National saving}\index{national saving} is output not spent on private or public consumption: $S = Y - C - G$. \emph{Private saving}\index{private saving}%
\phantomsection\label{def:g12:saving-and-investment:private-saving} is what households and firms keep from their income after taxes $T$ and consumption, $Y - T - C$.
\emph{Public saving}\index{public saving}\phantomsection\label{def:g12:saving-and-investment:public-saving} is the government's revenue minus its spending, $T - G$:
negative when it runs a deficit.
\end{definition}

\begin{proposition}[Saving and investment]\label{prop:g12:saving-and-investment:identity}
\omterm{def:g12:saving-and-investment:national-saving}{National saving} is the sum of private and \omterm{def:g12:saving-and-investment:national-saving}{public saving}. In a closed economy ($X = M = 0$), \omterm{def:g12:saving-and-investment:national-saving}{national saving} equals investment; in an open economy, $S - I = X - M$.
\end{proposition}
\begin{proof}
$(Y - T - C) + (T - G) = Y - C - G = S$. From $Y = C + I + G + X - M$, $S = Y - C - G = I + X - M$; with $X = M = 0$, $S = I$.
\end{proof}

\begin{definition}[Net lending]\label{def:g12:saving-and-investment:net-lending}
A sector's \emph{net lending}\index{net lending} is its saving minus its investment: positive when it has funds to lend to other sectors (a financing capacity), negative
when it must borrow (a financing need).
\end{definition}

\begin{measure}[Net lending in the accounts]
National accounts divide the economy into sectors: households, firms, financial institutions, government and the rest of the world. For each they record saving and
investment, and the difference, \omterm{def:g12:saving-and-investment:net-lending}{net lending} or net borrowing. Households are usually net lenders, firms net borrowers, the government a borrower when in deficit; the
sectors' \omterm{def:g12:saving-and-investment:net-lending}{net lending} adds up to the nation's, and that is lent to or borrowed from the rest of the world.
\end{measure}

\section{The market for loanable funds}\label{sec:g12:saving-and-investment:market}

\begin{definition}[Market for loanable funds]\label{def:g12:saving-and-investment:loanable-funds}
The \emph{market for loanable funds}\index{market for loanable funds} is the model of a market in which savers supply funds and investors demand them, at a price that is the
\omterm{def:g10:price-indices:real-interest-rate}{real interest rate}.
\end{definition}

Saving rises with the real rate, since each dollar saved earns more; investment falls with it, since fewer projects earn more than the cost of borrowing (the \omterm{def:g12:present-value:present-value}{present value} of
\cref{ch:g12:present-value}). The rate settles where they are equal.

\begin{definition}[Crowding out]\label{def:g12:saving-and-investment:crowding-out}
\emph{Crowding out}\index{crowding out} is the fall in private investment that follows a rise in government borrowing, because the borrowing raises the \omterm{def:g10:price-indices:real-interest-rate}{real interest rate}.
\end{definition}

\begin{omfigure}
\begin{tikzpicture}
\begin{axis}[width=0.66\linewidth, height=5.4cm, xmin=10, xmax=30, ymin=0, ymax=5, xlabel={funds, per cent of GDP}, ylabel={real interest rate, per cent},
  xtick={19,20}, ytick={2,2.5}, grid=major, grid style={gray!20}]
\addplot[omDef, thick, domain=0:5, samples=2] ({24 - 2*x}, {x}) node[pos=0.05, above right, font=\scriptsize] {investment};
\addplot[omThm, thick, domain=0:5, samples=2] ({16 + 2*x}, {x}) node[pos=0.9, right, font=\scriptsize] {saving};
\addplot[omThm, thick, dashed, domain=0:5, samples=2] ({14 + 2*x}, {x}) node[pos=0.9, left, font=\scriptsize] {saving with a deficit};
\fill (axis cs:20,2) circle (1.5pt) (axis cs:19,2.5) circle (1.5pt);
\end{axis}
\end{tikzpicture}

{\small A deficit of 2\,\% of GDP, with the numbers of the weekend problem. \omterm{def:g12:saving-and-investment:national-saving}{National saving} shifts left by the deficit; the rate rises from 2 to 2.5\,\%; \omterm{def:g12:saving-and-investment:national-saving}{private saving} rises
by 1 point, investment falls by 1 point. Half the deficit is financed by \omterm{def:g12:saving-and-investment:crowding-out}{crowding out} investment.}
\end{omfigure}

\begin{model}{The loanable-funds market}
\assumes A closed economy; saving rises and investment falls with the \omterm{def:g10:price-indices:real-interest-rate}{real interest rate}.
\predicts \omterm{def:g12:saving-and-investment:national-saving}{National saving} and investment are equal in each country; a government deficit raises the real rate and crowds out investment.
\tested In the world of 2010--2019 the closed-economy prediction fails: across 155 countries, investment rose by only about a quarter of a point for each extra point of
  \omterm{def:g12:saving-and-investment:national-saving}{national saving} (the opening chart). Feldstein and Horioka (1980) asked whether a country's saving stays at home or flows to wherever it earns most; for 16 rich countries
  in 1960--74 they found that investment rose almost point for point with saving, a coefficient of 0.89. Across 155 countries in 2010--19 the slope
  is only 0.25: a country's own saving accounts for a small part of the differences in investment.
% ledger: b1:feldstein1980
\breaks Where capital flows freely, a country can borrow abroad at the world rate, so a deficit raises borrowing from abroad more than the rate; and when the rate is near
  zero and cannot fall, the market does not bring saving and investment together through the rate.
\end{model}
% ledger: g12:fh1980

\section{How firms pay for investment}\label{sec:g12:saving-and-investment:firms}

\begin{definition}[Gross operating surplus and self-financing]\label{def:g12:saving-and-investment:gross-operating-surplus}
A firm's \emph{gross operating surplus}\index{gross operating surplus} is its value added minus the wages and taxes on production it pays: what is left to pay interest,
profits and taxes on profits, and to invest. \emph{Self-financing}\index{self-financing}\phantomsection\label{def:g12:saving-and-investment:self-financing} is investment
paid for out of the firm's own saving, the profits it keeps.
\end{definition}

\begin{definition}[Financial intermediary and share]\label{def:g12:saving-and-investment:intermediary}
A \emph{financial intermediary}\index{financial intermediary} is an institution, such as a bank, an insurer or a fund, that collects savings and lends or invests them. A
\emph{share}\index{share (in a company)}\phantomsection\label{def:g12:saving-and-investment:share} in a company is a part of its ownership, which gives a claim on its
profits and a vote on its management.
\end{definition}

A firm that needs more than it keeps can borrow from a bank, sell \omterm{def:g12:present-value:bond}{bonds}, or sell new \omterm{def:g12:saving-and-investment:intermediary}{shares}. A loan or a \omterm{def:g12:present-value:bond}{bond} must be repaid with interest whatever happens; a \omterm{def:g12:saving-and-investment:intermediary}{share} pays
dividends only if there are profits, so it shares the risk with the investor. Banks lend to small firms whose quality only a close lender can judge; large, well-known firms
can borrow directly from savers by selling \omterm{def:g12:present-value:bond}{bonds} and \omterm{def:g12:saving-and-investment:intermediary}{shares}.

\begin{omfigure}
\begin{tikzpicture}
\begin{axis}[xbar stacked, /pgf/bar width=12pt, width=0.7\linewidth, height=3.6cm, xmin=0, xmax=100, xlabel={per cent of investment},
  ytick={1,2}, yticklabels={{young firm}, {mature firm}}, enlarge y limits=0.5, legend style={at={(0.5,1.05)}, anchor=south, font=\scriptsize, draw=none},
  legend columns=4]
\addplot[fill=omDef!60, draw=omDef] coordinates {(20,1) (70,2)};
\addplot[fill=omThm!60, draw=omThm] coordinates {(50,1) (10,2)};
\addplot[fill=omProp!60, draw=omProp] coordinates {(0,1) (15,2)};
\addplot[fill=omMeth!60, draw=omMeth] coordinates {(30,1) (5,2)};
\legend{own saving, bank loans, bonds, new shares}
\end{axis}
\end{tikzpicture}

{\small How two invented firms pay for a year's investment. The young firm keeps little profit and depends on its bank and on new owners; the mature firm pays for most of
its investment itself and borrows the rest by selling \omterm{def:g12:present-value:bond}{bonds}.}
\end{omfigure}

\section{When finance fails}\label{sec:g12:saving-and-investment:crises}

\begin{definition}[Speculative bubble]\label{def:g12:saving-and-investment:bubble}
A \emph{speculative bubble}\index{speculative bubble} is a rise in an asset's price above the \omterm{def:g12:present-value:present-value}{present value} of the income it will bring, sustained by buyers' expectation that
the price will keep rising.
\end{definition}

A fall in asset prices reaches output and jobs through several channels. Households who feel poorer spend less. Borrowers whose collateral loses value can borrow less. Banks
that lose on loans lend less to everyone, and lenders that must sell assets quickly push prices down further, hurting other holders.

\begin{omfigure}
\begin{tikzpicture}[font=\scriptsize, node distance=0.5cm, box/.style={draw, rounded corners, fill=omDef!10, align=center, minimum height=0.9cm, text width=2.0cm}]
\node[box] (a) {asset prices fall};
\node[box, right=of a] (b) {wealth and collateral fall};
\node[box, right=of b] (c) {losses at banks; fire sales};
\node[box, right=of c] (d) {less credit, less spending};
\node[box, below=0.6cm of d, fill=omThm!10] (e) {output and jobs fall};
\draw[->, thick] (a) -- (b); \draw[->, thick] (b) -- (c); \draw[->, thick] (c) -- (d); \draw[->, thick] (d) -- (e);
\draw[->, thick, dashed] (e.west) -| (a.south) node[pos=0.25, below, font=\scriptsize] {more defaults};
\end{tikzpicture}

{\small How a fall in asset prices reaches the real economy, and feeds back on itself.}
\end{omfigure}

\begin{omchart}{United States: real GDP per person (index, 100 in the year before the crisis, solid) and the unemployment rate (per cent, dashed), 1926--1940 and
2004--2016.}{Jordà--Schularick--Taylor Macrohistory Database}
\begin{tikzpicture}
\begin{groupplot}[group style={group size=2 by 1, horizontal sep=1.3cm}, omtime, height=5.0cm, ymin=0, ymax=110]
\nextgroupplot[omtime, width=0.48\linewidth, xmin=1926, xmax=1940, xtick={1926,1930,...,1938}, ylabel={index; per cent}]
\addplot[omDef, thick] table[x=year, y=gdp] {figdata/grade-12/07-saving-and-investment/out/crisis1929.csv};
\addplot[omThm, thick, dashed] table[x=year, y=unemp] {figdata/grade-12/07-saving-and-investment/out/crisis1929.csv};
\nextgroupplot[omtime, width=0.48\linewidth, xmin=2004, xmax=2016, xtick={2004,2008,2012,2016}]
\addplot[omDef, thick] table[x=year, y=gdp] {figdata/grade-12/07-saving-and-investment/out/crisis2007.csv} node[pos=0.5, below, font=\scriptsize] {GDP per person};
\addplot[omThm, thick, dashed] table[x=year, y=unemp] {figdata/grade-12/07-saving-and-investment/out/crisis2007.csv} node[pos=0.5, above, font=\scriptsize] {unemployment};
\end{groupplot}
\end{tikzpicture}
\end{omchart}
% figdata: crisis_key (gdp1933 69.2, u1929 3.2, u1933 24.9, gdp2009 94.5, u2007 4.6, u2010 9.6, house2012 84.5)

\begin{record}[1929--33 and 2007--09: from a crash to output and jobs]
% ledger: g12:crash1929, g12:gr2007, g12:lehman
The Dow Jones index of American \omterm{def:g12:saving-and-investment:intermediary}{shares} closed at 381 on 3 September 1929; it fell by nearly 13\,\% on 28 October and by nearly 12\,\% the next day, and by 8 July 1932 it stood at
41, 89\,\% below its peak. Real output per person fell by almost a third from 1929 to 1933, and unemployment rose from 3.2\,\% to 24.9\,\%. In the second crisis, American
home prices fell by roughly 30\,\% from their peak in mid-2006 to mid-2009; the investment bank Lehman Brothers filed for bankruptcy on 15 September 2008, the largest
filing in American history. The recession ran from December 2007 to June 2009; real GDP fell by 4.3\,\% and unemployment rose from 5\,\% to a peak of 10\,\% in October
2009. The second crisis was smaller than the first; the chart shows how much.
\end{record}

\section{Exercises}

\begin{exercise}[$\star$]\label{exo:g12:saving-and-investment:1}
A closed economy has $Y = 1000$, $C = 600$ and $G = 200$. Find \omterm{def:g12:saving-and-investment:national-saving}{national saving} and investment.
\end{exercise}

\begin{exercise}[$\star$]\label{exo:g12:saving-and-investment:2}
An open economy has $Y = 1000$, $C = 650$, $G = 200$ and $I = 180$. Find \omterm{def:g12:saving-and-investment:national-saving}{national saving} and \omterm{def:g10:national-accounts:net-exports}{net exports}.
\end{exercise}

\begin{exercise}[$\star$]\label{exo:g12:saving-and-investment:3}
From the opening chart, which country saves much more than it invests, and which invests much more than it saves? What does each do with the difference?
\end{exercise}

\begin{exercise}[$\star$]\label{exo:g12:saving-and-investment:4}
$Y = 1000$, $T = 250$, $C = 600$, $G = 280$. Find private, public and \omterm{def:g12:saving-and-investment:national-saving}{national saving}.
\end{exercise}

\begin{exercise}[$\star$]\label{exo:g12:saving-and-investment:5}
Households save 120 and invest 40 (in housing); firms save 90 and invest 150. Find each sector's \omterm{def:g12:saving-and-investment:net-lending}{net lending}, and their total.
\end{exercise}

\begin{exercise}[$\star\star$]\label{exo:g12:saving-and-investment:6}
Saving is $S = 100 + 20r$ and investment $I = 180 - 20r$, $r$ in per cent. Find the equilibrium rate and amount.
\end{exercise}

\begin{exercise}[$\star\star$]\label{exo:g12:saving-and-investment:7}
In \cref{exo:g12:saving-and-investment:6}, firms become more optimistic and investment becomes $I = 200 - 20r$. Find the new equilibrium.
\end{exercise}

\begin{exercise}[$\star\star$]\label{exo:g12:saving-and-investment:8}
In \cref{exo:g12:saving-and-investment:6}, the government runs a deficit of 20. Find the new rate, investment and \omterm{def:g12:saving-and-investment:national-saving}{private saving}, and the \omterm{def:g12:saving-and-investment:intermediary}{share} of the deficit crowded out.
\end{exercise}

\begin{exercise}[$\star\star$]\label{exo:g12:saving-and-investment:9}
In \cref{exo:g12:saving-and-investment:6}, households decide to save 40 more at every rate. Find the new rate and investment. What would happen in an open economy at a fixed
world rate?
\end{exercise}

\begin{exercise}[$\star\star$]\label{exo:g12:saving-and-investment:10}
A firm invests 500, keeps profits of 350 and borrows the rest. Find its \omterm{def:g12:saving-and-investment:gross-operating-surplus}{self-financing} ratio.
\end{exercise}

\begin{exercise}[$\star\star$]\label{exo:g12:saving-and-investment:11}
From the chart of the two crises, give the index of GDP per person and the unemployment rate in 1933 and in 2009, and say which crisis was deeper.
\end{exercise}

\begin{exercise}[$\star\star\star$]\label{exo:g12:saving-and-investment:12}
What evidence contradicts the closed-economy model of loanable funds? Use the opening chart and its fitted slope.
\end{exercise}

\begin{exercise}[$\star\star\star$]\label{exo:g12:saving-and-investment:13}
An asset pays 5 a year for ever, and the interest rate is 5\,\%. Its price is 150, and buyers expect it to rise by 10\,\% a year. What is its \omterm{def:g12:present-value:present-value}{present value}? Show that a buyer
earns more than 5\,\% as long as the expectation holds. What happens when buyers stop expecting the rise?
\end{exercise}

\begin{exercise}[$\star\star\star$]\label{exo:g12:saving-and-investment:14}
A bank holds assets of 100 and capital of 8. It must sell 30 of assets quickly, and fire sales fetch 80\,\% of their value. What does it lose, and what is left of its
capital? What happens to other banks that hold the same assets?
\end{exercise}

\begin{exercise}[$\star\star\star$]\label{exo:g12:saving-and-investment:15}
In the open economy of the weekend problem, why does a deficit leave investment unchanged? Who lends the government the money, and what does the country owe in return?
\end{exercise}

\section{Problem: Brenland's Loanable Funds}

\begin{problem}[{Weekend problem --- a deficit, closed and open}]
\label{pb:g12:saving-and-investment:1}
In Brenland, \omterm{def:g12:saving-and-investment:national-saving}{private saving} is $S_p = 16 + 2r$ and investment $I = 24 - 2r$, both in per cent of GDP, with $r$ the \omterm{def:g10:price-indices:real-interest-rate}{real interest rate} in per cent. Brenland's GDP is 393\,264
million dalens.

\medskip
\noindent\textbf{Part I --- A balanced budget, closed economy.}
\begin{enumerate}
\item With a balanced budget, find the equilibrium real rate.
\item Find saving and investment, in per cent of GDP and in dalens.
\item What is \omterm{def:g12:saving-and-investment:national-saving}{public saving}? \omterm{def:g12:saving-and-investment:national-saving}{National saving}?
\item Firms invest more if the rate falls. Why, using \omterm{def:g12:present-value:present-value}{present value}?
\item Households save more if the rate rises. Give one reason.
\end{enumerate}

\medskip
\noindent\textbf{Part II --- A deficit, closed economy.}
\begin{enumerate}[resume]
\item The government runs a deficit of 2\,\% of GDP. Write \omterm{def:g12:saving-and-investment:national-saving}{national saving} as a function of $r$.
\item Find the new equilibrium rate.
\item Find investment and \omterm{def:g12:saving-and-investment:national-saving}{private saving}.
\item By how much has investment fallen? What is the crowding-out ratio, investment lost per unit of deficit?
\item Who finances the rest of the deficit?
\end{enumerate}

\medskip
\noindent\textbf{Part III --- Open to Kestria.}
\begin{enumerate}[resume]
\item Capital now flows freely between Brenland and Kestria, and Brenland is small enough to borrow at a world real rate of 2\,\%. With the deficit, find investment and
  \omterm{def:g12:saving-and-investment:national-saving}{national saving}.
\item Find $S - I$, and so \omterm{def:g10:national-accounts:net-exports}{net exports}, in per cent of GDP and in dalens.
\item Who lends Brenland the difference?
\item Compare the effects of the deficit on investment in the closed and the open economy.
\item What does Brenland owe in future years as a result?
\end{enumerate}

\medskip
\noindent\textbf{Part IV --- Interpretation.}
\begin{enumerate}[resume]
\item The opening chart shows saving and investment loosely related across countries. Which part of the problem does that support?
\item If the world rate rose to 3\,\%, what would Brenland's investment be?
\item Why might a large country's deficit raise the world rate?
\item Under what condition would the deficit pay for itself by raising future output?
\item State the result of the problem: the crowding-out ratio in the closed economy, and the deficit of \omterm{def:g10:national-accounts:net-exports}{net exports} once capital flows.
\end{enumerate}
\end{problem}
